% =============================================================================
% architectural_form_section.tex
% Draft for §4 (The Architectural Form) of paper-reordered.tex.
%
% This is a SEPARATE working file. paper-reordered.tex is NOT modified.
%
% =============================================================================
% SCOPE OF THIS DRAFT
% =============================================================================
% All three main components are now developed:
%   §4.1 Experimental Control            (Orchestration, Stimulus)
%   §4.2 Temporal Observation            (Acquisition, Coordination)
%   §4.3 Representation and Artifact     (Representation, Analysis, Retention)
%        Lifecycle
%
% Update log:
%   - 2026-05-20 (a): drafted Component 3 only.
%   - 2026-05-20 (b): extended with Components 1 and 2, Mapping subsection
%     filled, §3.1 coverage re-verified at the full-form level.
%
% =============================================================================
% LOCKED DECISIONS (this session)
% =============================================================================
%   - Hierarchy: 3 main components.
%       §4.1  Experimental Control            (Orchestration, Stimulus)         [DRAFTED]
%       §4.2  Temporal Observation            (Acquisition, Coordination)       [DRAFTED]
%       §4.3  Representation and Artifact     (Representation, Analysis,        [DRAFTED]
%             Lifecycle                        Retention)
%   - Depth: per sub-component blocks. Each sub-component carries its own
%     Purpose / Responsibilities / Invariant.
%   - Bullet block label: "Responsibilities" (was "absorbs" in last turn's
%     analysis — same items, different word).
%   - When integrated into paper-reordered.tex, Component 3 becomes §4.3.
%     Until Components 1 and 2 are drafted, the \subsection here will render as
%     §4.1; that is expected and harmless during drafting.
%
% =============================================================================
% PROVENANCE
% =============================================================================
%   - Items inside each Responsibilities block are pulled from the structural
%     paper (A_Reproducible_Architecture_..._FORM.pdf, §III).
%   - Each item carries a comment naming its origin in §III.{layer}.{item}.
%   - Purpose / Invariant text is drafted from the last-turn analysis and is
%     expected to be refined.
%
% =============================================================================
% FILE LAYOUT
% =============================================================================
%   PART A — Skeleton (placeholders, no content).
%   PART B — Filled skeleton (placeholders replaced with items from §III.E/F/G).
%   NOTES  — §3 mapping cross-check, open questions, things to flag in next pass.
% =============================================================================



% =============================================================================
% PART A — SKELETON
% =============================================================================

\section{The Architectural Form}
% ----- Opening sentence (1) -----
% TODO: one sentence naming what this section delivers (the form itself).

% ----- Explanation of what we mean by "form" (1.1) -----
% TODO: 2-3 sentences clarifying that the form is not a system blueprint
% but a specification of roles, sub-components, and invariants that any
% realization must satisfy.

% ----- Roadmap of subsections (1.2) -----
% TODO: 1 sentence naming the three main components and signalling the
% order in which the section traverses them. Note: this draft only develops
% Component 3.

\subsection{Experimental Control}
% Component 1 opening: 2-3 sentences naming the role and the two
% sub-components it decomposes into.
% TODO

  \subsubsection{Orchestration}
  \textbf{Purpose.} % TODO: one sentence — what this sub-component does.

  \textbf{Responsibilities.}
  \begin{itemize}
    \item % TODO: pull item from §III.A.1
    \item % TODO: pull item from §III.A.2
    \item % TODO: pull item from §III.A.3
  \end{itemize}

  \textbf{Invariant.} % TODO: one sentence — the property realizations must preserve.

  \textbf{Cross-reference.} % TODO: link to §4.1.2 (Stimulus) where relevant.

  \subsubsection{Stimulus}
  \textbf{Purpose.} % TODO

  \textbf{Responsibilities.}
  \begin{itemize}
    \item % TODO: pull item from §III.B (deduplicated; drop "Temporal alignment of actions" which belongs in §4.1.1)
    \item % TODO
    \item % TODO
  \end{itemize}

  \textbf{Invariant.} % TODO

  \textbf{Cross-reference.} % TODO: link to §4.1.1 (Orchestration) where relevant.

\subsection{Temporal Observation}
% Component 2 opening: 2-3 sentences naming the role and the two
% sub-components it decomposes into.
% TODO

  \subsubsection{Acquisition}
  \textbf{Purpose.} % TODO

  \textbf{Responsibilities.}
  \begin{itemize}
    \item % TODO: pull item from §III.C.1
    \item % TODO: pull item from §III.C.2
    \item % TODO: pull item from §III.C.3
  \end{itemize}

  \textbf{Invariant.} % TODO

  \textbf{Cross-reference.} % TODO: link to §4.2.2 (Coordination) where relevant.

  \subsubsection{Coordination}
  \textbf{Purpose.} % TODO

  \textbf{Responsibilities.}
  \begin{itemize}
    \item % TODO: pull item from §III.D.1
    \item % TODO: pull item from §III.D.2
    \item % TODO: pull item from §III.D.3
  \end{itemize}

  \textbf{Invariant.} % TODO

  \textbf{Cross-reference.} % TODO: link to §4.2.1 (Acquisition) where relevant.

\subsection{Representation and Artifact Lifecycle}
% Component 3 opening: 2-3 sentences naming the role and the three
% sub-components it decomposes into.
% TODO

  \subsubsection{Representation}
  \textbf{Purpose.} % TODO: one sentence — what this sub-component does.

  \textbf{Responsibilities.}
  \begin{itemize}
    \item % TODO: pull item from §III.E.1
    \item % TODO: pull item from §III.E.2
    \item % TODO: pull item from §III.E.3
  \end{itemize}

  \textbf{Invariant.} % TODO: one sentence — the property realizations must preserve.

  \subsubsection{Analysis}
  \textbf{Purpose.} % TODO

  \textbf{Responsibilities.}
  \begin{itemize}
    \item % TODO: pull item from §III.F.1
    \item % TODO: pull item from §III.F.2
    \item % TODO: pull item from §III.F.3
  \end{itemize}

  \textbf{Invariant.} % TODO

  \subsubsection{Retention}
  \textbf{Purpose.} % TODO

  \textbf{Responsibilities.}
  \begin{itemize}
    \item % TODO: pull item from §III.G.1
    \item % TODO: pull item from §III.G.2
    \item % TODO: pull item from §III.G.3
  \end{itemize}

  \textbf{Invariant.} % TODO

\subsection*{Mapping to Requirements (§3)}
% TODO: for each sub-component, cite which §3 requirements its Invariant
% satisfies. Decide whether this subsection lives inside §4.3 or moves up
% to §4 once Components 1 and 2 are drafted.



% =============================================================================
% PART B — FILLED SKELETON
% =============================================================================
% Same structure as Part A; placeholders replaced with first-pass text.
% Purpose / Invariant text is drafted and expected to be refined.
% Items in Responsibilities are verbatim from the structural paper §III.E/F/G,
% with the source flagged in a trailing comment.
% =============================================================================

\section{The Architectural Form}

% ----- Opening sentence (1) -----
% [first-pass — refine later]
The architectural form for repeated volatile-memory observation is specified as a set of essential roles, each defined by its responsibilities and by the invariants that any valid realization must preserve.

% ----- Explanation of what we mean by "form" (1.1) -----
% [first-pass — refine later]
By \emph{form} we mean a specification at the level of architectural roles and relations rather than at the level of mechanisms or implementations. The form does not prescribe how each role is realized; it identifies what roles must exist, how they decompose into sub-components, and what properties must remain invariant across realizations.

% ----- Roadmap of subsections (1.2) -----
% [first-pass — refine later]
The form decomposes into three main components: \emph{Experimental Control}, which produces the conditions under which observation is meaningful; \emph{Temporal Observation}, which samples and conveys system states without distorting their temporal relationships; and \emph{Representation and Artifact Lifecycle}, which transforms observed sequences into analysis-ready intermediate forms and governs their persistence. The remainder of this section develops each in turn, and concludes with a mapping from the form's invariants back to the requirements of \S3.

\subsection{Experimental Control}
% [Opener polished 2026-05-20:
%   - "Component 1" subject dropped (user preference); subject becomes
%     the component name itself, "Experimental Control".
%   - "produces" replaced with "establishes" (structural verb,
%     consistent with sub-components below).
%   - Artifacts gloss added: step boundaries (Orchestration) and
%     workload specifications (Stimulus) introduced as defined
%     concepts before sub-components reference them.
%   - Italicized on first use: roles (Orchestration, Stimulus) and
%     artifacts (step boundaries, workload specifications).
%   - Compressed role-description claims removed from opener; they
%     are carried by the sub-components below.]
Experimental Control establishes the controlled conditions under which observation becomes temporally and contextually meaningful. It decomposes into two sub-components, each providing a distinct structural artifact: \emph{Orchestration} provides the \emph{step boundaries} that mark the time intervals of the experiment, and \emph{Stimulus} provides the \emph{workload specifications} that describe what runs inside each step.

  \subsubsection{Orchestration}
  % [Purpose polished 2026-05-20 in two rounds:
  %   R1: expanded from one sentence to a 3-sentence claim-to-consequence
  %       paragraph. User accepted side-by-side.
  %   R2: regression fix — "the phase in which it was produced" changed to
  %       "the phase under which it was observed". The causal verb "produced"
  %       was the same pattern we removed from the Responsibilities;
  %       structural-observation phrasing keeps the semantic-free stance.]
  \textbf{Purpose.} The \emph{Orchestration} role imposes a stable ordering on experimental actions and preserves the boundaries between successive experimental phases. By keeping the ordering and boundaries of execution under explicit control rather than letting them emerge from downstream timing, this role ensures that each observed transition remains aligned with the phase under which it was observed. Consequently, later stages of the form can interpret a sequence of states as part of a defined progression rather than as an undifferentiated stream of execution.

  % [Responsibilities polished 2026-05-20 in three rounds:
  %   R1: short labels expanded into prose using structural paper "Polished
  %       paper prose" entries with prescriptive register stripped.
  %   R2: nouns concretized (workloads, memory states, steps, sampling)
  %       so reader has anchor points.
  %   R3: causal verbs (produced, caused) replaced with positional verbs
  %       (located, situated, recorded, aligned) to prevent the reviewer
  %       concern about semantic-free claims. Explicit structural-not-
  %       semantic stance kept in Bullet 2.]
  \textbf{Responsibilities.}
  \begin{itemize}
    \item \textbf{Controlled experiment sequencing.} The order in which experimental workloads are executed and memory states are sampled is determined by the Orchestration role itself, not by downstream processing. As a result, each observed state transition can be located within the controlled execution timeline rather than appearing as an isolated change.
      % [src: structural paper §III.A.1 — Orchestration Layer]
    \item \textbf{Boundary preservation between steps.} The start and end of each experimental step are recorded and propagated through the form, so that any observed transition can be situated within a specific step rather than within an unlabelled segment of execution. Step markers are structural; they identify where in the schedule a transition occurred, not what activity produced it.
      % [src: structural paper §III.A.2]
    \item \textbf{Temporal alignment of actions.} Control decisions, workload activity, and memory-state sampling are scheduled so that their timing relationships remain fixed and known. Without that stability, the observation sequence could not be aligned with the control timeline, and the structural association between transitions and steps would be lost.
      % [src: structural paper §III.A.3]
      % NOTE: in the source PDF this item also appeared duplicated under
      % §III.B.4 (Stimulus). Kept here only.
  \end{itemize}

  % [Invariant polished 2026-05-20 in two rounds:
  %   R1: subject swapped from "control" to "Orchestration"; verb tightened
  %       from "determined by" to "imposed by"; minor nominalization fixes.
  %   R2: clarity enhancement — concrete subjects (workloads, steps),
  %       downstream stages enumerated, failure mode split from practical
  %       consequence, neutral verbs (yield, align, shift, track) to
  %       preserve semantic-free stance. Audit checklist appended.]
  \textbf{Invariant.} The order in which experimental workloads run and the boundaries between successive steps are imposed by Orchestration, never by the timing of downstream stages such as transformation, analysis, or artifact handling. If downstream processing were allowed to shape this order, the temporal positions of transitions in the observation sequence would shift with downstream processing speed rather than tracking the controlled schedule. Two runs of the same experiment could then yield observation sequences that no longer align in time, breaking comparability between them. A realization satisfies this invariant when transformation latency, output generation, and artifact handling cannot change which memory state is sampled next, when the next experimental step begins, or where a step boundary is placed in time.
  % [working note — \S3 requirements this Invariant addresses (precise refs
  %  kept here for audit; user-facing Cross-reference uses light-formal-proof
  %  prose without these pointers):
  %   - \S3.1, P1: preserve phase and condition for each sampled state.
  %   - \S3.1, P3, S1: support distinguishable conditions and units.
  %   - \S3.1, P3, S3: preserve segments without semantic labels.
  %   - \S3.3: form-level separation of control from downstream roles.]

  % [Cross-reference polished 2026-05-20 in three rounds:
  %   R1: tightened structural link and named \S3.1 requirement.
  %   R2: expanded into bulleted requirement decomposition with explicit
  %       paragraph/sentence citations.
  %   R3: user wanted a light-formal-proof reading instead of citation;
  %       claims italicized, composition reasoning per claim,
  %       paragraph/sentence numbers removed from prose.]
  \textbf{Cross-reference.} The step boundaries established by Orchestration mark the time intervals within which Stimulus (\S4.1.2) places its workload definitions. To map this onto \S3.1's vocabulary: a phase is an experimental step (the time slice that Orchestration marks), and a condition is a workload definition (the activity that Stimulus places inside that slice). With that mapping in place, the two sub-components jointly satisfy three \S3.1 requirements.

  \emph{Preservation of phase and condition for each sampled state.} \S3.1 requires that each sampled state be linked to the phase it belongs to and the condition active during it. Orchestration provides the phase link through its step boundaries; Stimulus provides the condition link through its workload definitions. Any state sampled within step $s$ is therefore associated with both, satisfying the requirement.

  \emph{Distinguishability of conditions and units.} \S3.1 requires both distinguishable conditions (different workloads) and distinguishable units (different steps). Orchestration distinguishes units by emitting distinct step boundaries; Stimulus distinguishes conditions by attaching a distinct workload definition to each step. Both kinds of difference are preserved: differences between steps are marked by Orchestration, differences between workloads are marked by Stimulus.

  \emph{Preservation of segments without semantic labels.} \S3.1 requires that distinguishable segments be preserved without attaching meaning to them. Orchestration's step markers are structural (a step number, not a description of what happens), and Stimulus's workload definitions are structural specifications rather than interpretations. The segments are therefore preserved using only structural markers, and no semantic label is introduced.

  \subsubsection{Stimulus}
  % [Purpose polished 2026-05-20 in two rounds:
  %   R1: expanded to claim-to-consequence shape matching §4.1.1 Orchestration.
  %   R2: concrete glosses added on first use — workload = (program, inputs,
  %       parameters); schedule = (order and timing). Self-containment
  %       improved per user request. Stance-safe: "what a workload does"
  %       avoided in favor of structural specification.]
  \textbf{Purpose.} The \emph{Stimulus} role introduces controlled, reproducible workload definitions into the system, where each workload runs within a step boundary supplied by Orchestration. By keeping the workload (the program, inputs, and parameters being run) separate from the schedule (the order and timing of when each workload runs), this role ensures the two remain independent: the workload can be changed without redesigning the schedule, and the schedule can be changed without redefining any workload. Consequently, a single schedule can host many different workloads, and a single workload can be reused across many schedules.

  % [Responsibilities polished 2026-05-20 in four rounds:
  %   R1: short labels expanded into prose using structural paper "Polished
  %       paper prose" entries.
  %   R2: nouns concretized; experimenter agency named.
  %   R3: removed "experimenter" (form-neutral about who fixes workloads);
  %       removed overclaim about identical memory evolution across runs
  %       (volatile memory varies); Bullet 3 cleaned with parallel "decides"
  %       verbs and parenthetical structural outputs.
  %   R4: Bullet 2 trailing phrase swapped to positive-consequence ending
  %       (Option 2) so the bullet states what reproducibility delivers
  %       at form level without referencing memory or any hypothetical
  %       experiment.]
  \textbf{Responsibilities.}
  \begin{itemize}
    \item \textbf{Controlled behavioral inputs.} The activity introduced into the system is bounded to a defined set of workloads that are fixed before any run begins. Because the workload set is known in advance of any memory state being sampled, the conditions under which states evolve can be compared across runs of the same experiment.
      % [src: structural paper §III.B.5 — Stimulus Layer; first non-duplicate item]
    \item \textbf{Reproducible workload definition.} Each workload is described by an explicit specification (program, inputs, parameters, and any setup it requires), and that specification is preserved across runs. The form's reproducibility property covers the workload specification and its invocation at fixed step positions, so the same workload definition enters the system at the same step each time.
      % [src: structural paper §III.B.6]
    \item \textbf{Separation of stimulus from control.} The Stimulus role decides what runs in the system; the Orchestration role decides when each workload runs. Each role provides its own structural output (Stimulus the workload specifications, Orchestration the step boundaries), and neither role provides the other's output. As a result, changes to a workload do not require changes to the schedule, and changes to the schedule do not require changes to any workload.
      % [src: structural paper §III.B.7]
  \end{itemize}

  % [Invariant polished 2026-05-20 in two rounds:
  %   R1: subject swapped to "every workload running"; forbidden direction
  %       named explicitly; failure-mode and audit-target sentences added.
  %   R2: agent-fix — replaced "If the schedule were allowed to alter..."
  %       with "If the orchestration mechanism were allowed to alter...".
  %       The schedule is a passive artifact, not an agent; the orchestration
  %       mechanism is the role that could act. Keeps Invariant internally
  %       consistent (same agent named in sentence 1).]
  \textbf{Invariant.} Every workload running in the system is defined by an explicit specification provided by Stimulus, never derived from or modified by the orchestration mechanism. If the orchestration mechanism were allowed to alter the workload definition mid-experiment (for example, in response to observed memory state or downstream feedback), the workload would cease to be a fixed input and would instead become a function of the schedule, breaking the separation between what runs and when it runs. A realization satisfies this invariant when the orchestration mechanism cannot read, modify, or replace any workload specification, and when no workload specification can read or alter the schedule of step boundaries.
  % [working note — \S3 requirements this Invariant addresses:
  %   - \S3.1, P3, S2: differences in memory evolution should be relatable
  %     to differences in observed system activity without semantic knowledge.
  %   - \S3.3: separation of stimulus from control is itself one of the
  %     form-level separations that must not collapse.]

  % [Cross-reference polished 2026-05-20 in three rounds:
  %   R1: mirrored §4.1.1 light-formal-proof structure; back-references
  %       §4.1.1 for joint §3.1 proof to avoid duplication; new §3.3
  %       separation claim block added.
  %   R2: "this sub-component" replaced with "the following" per user;
  %       "emit" replaced with "provide" for consistency with Responsibilities
  %       Bullet 3; final sentence restructured into two-level summary
  %       (role decomposition + Invariant).
  %   R3: "collapse into a single undifferentiated process" replaced with
  %       plain "merged into one combined role" for clarity.]
  \textbf{Cross-reference.} Stimulus's workload definitions are placed inside the step boundaries supplied by Orchestration (\S4.1.1). The joint satisfaction of \S3.1's requirements by Orchestration and Stimulus is shown at \S4.1.1; the following additionally satisfies the form-level separation requirement of \S3.3.

  \emph{Preservation of the stimulus-versus-control separation.} \S3.3 requires that distinct architectural roles not be merged into one combined role, since each role protects different responsibilities. Stimulus and Orchestration each provide their own structural output: Stimulus provides workload specifications, Orchestration provides step boundaries. Neither role can substitute for the other, and the Invariant above forbids the orchestration mechanism from modifying any workload specification. The separation is therefore preserved at two levels: by the form's role decomposition (Stimulus and Orchestration are distinct sub-components), and by the Invariant above (no role may alter another's output).

\subsection{Temporal Observation}
% Component 2 opening.
% [first-pass — refine later]
Component~2 samples system states over time and conveys them through the form while preserving their temporal relationships and observational context. It decomposes into two sub-components: \emph{Acquisition}, which captures successive states under a defined sampling regime that preserves adjacency and remains neutral to downstream use; and \emph{Coordination}, which passes captured states between stages whose processing rates may differ, without losing their order or context.

  \subsubsection{Acquisition}
  % [Purpose polished 2026-05-20 in three rounds:
  %   R1: expanded to claim-to-consequence shape matching §4.1.1 / §4.1.2;
  %       introduced observation sequence as a defined artifact on first use.
  %   R2: consequence framing changed from "downstream methods can consume
  %       without influencing the regime" to "the sequence is stable in shape
  %       regardless of which downstream methods consume it" (Variant B).
  %   R3: "regime" replaced with "protocol" (more specified, matches the
  %       fixed-in-advance constraint); "adjacency" kept and glossed
  %       concretely on first use (the structural paper's term, and the
  %       Cross-reference downstream also uses it).]
  \textbf{Purpose.} The \emph{Acquisition} role captures successive memory states over time and emits an \emph{observation sequence} in which the adjacency of consecutive states (their position next to each other in time) is preserved. By performing this capture under a sampling protocol that is fixed in advance of any downstream stage, this role keeps the temporal structure of the sequence determined entirely by the protocol, not by downstream processing. The sequence is therefore stable in shape regardless of which downstream methods consume it.

  % [Responsibilities polished 2026-05-20 in three rounds:
  %   R1: short labels expanded into bold-lead-in prose using structural
  %       paper "Polished paper prose" entries, with structural verbs only.
  %   R2: Bullet 1 framing made two-sided — contrasts against both the
  %       single-snapshot and the loose-collection extremes. Adds the
  %       "structured sequence" hint (known temporal position + explicit
  %       neighbour relationship) that points toward the uniformly-sampled-
  %       signal frame without committing the form to uniform-only.
  %   R3: Bullet 1 polish — "captured" → "sampled" (matches bullet label);
  %       compound clause consolidated into a relative clause on protocol;
  %       "operate as a whole" → "operate on the sequence rather than on
  %       individual samples"; contrast clause uses commas instead of
  %       parentheses (Option B).
  %   R4: Bullet 1 register softening — "loose collection of captures" →
  %       "unstructured collection of captures without explicit temporal
  %       relationships". "Loose" carried a faint evaluative charge that
  %       might read as aggressive toward snapshot-based approaches;
  %       "unstructured + without explicit temporal relationships" keeps
  %       the contrast two-sided while staying purely descriptive.]
  \textbf{Responsibilities.}
  \begin{itemize}
    \item \textbf{Repeated state sampling.} Memory states are sampled at successive points in time according to a sampling protocol that assigns each sample a known temporal position and an explicit relationship to its neighbours. The result is a structured sequence of samples, neither a single snapshot nor an unstructured collection of captures without explicit temporal relationships; subsequent stages of the form operate on the sequence rather than on individual samples.
      % [src: structural paper §III.C.1 — Acquisition Layer]
    % [Bullet 2 polished 2026-05-20 in multiple rounds:
    %   R1: short label expanded into "Consecutive samples ... interpolate, drop,
    %       or reorder ... bag of unrelated observations".
    %   R2: rewritten in Bullet 1's style (structural identity claim +
    %       downstream consequence with semicolon).
    %   R3: scope built into the claim itself ("At Acquisition's output, ...")
    %       to avoid the defensive "downstream may transform" disclaimer that
    %       sounded like the form was giving up on adjacency at later stages.
    %       The form-vs-realization boundary is implicit in the scope marker;
    %       downstream stages (Analysis etc.) have their own invariants.
    %   R4: ending swapped from "memory evolution" to "order and temporal
    %       placement of the captured states" — purely structural, avoids the
    %       semantic-loading risk that "evolution" carries.]
    \item \textbf{Preservation of state adjacency.} At Acquisition's output, the order of samples must match the order in which they were captured, with no interpolation, dropping, or reordering of samples. The result is an adjacency that subsequent stages of the form can rely on as a faithful record of the order and temporal placement of the captured states.
      % [src: structural paper §III.C.2]
    % [Bullet 3 polished 2026-05-20 in multiple rounds:
    %   R1: drafted in Bullet 1's style with scope marker "At Acquisition's
    %       output" and three forbidden operations (interpretation,
    %       transformation, labeling) symmetrically permitted downstream.
    %   R2: scope-only redraft — dropped the "subsequent stages may operate
    %       independently" disclaimer that read as defensive / handing off
    %       responsibility. Replaced with a positive characterization of
    %       Acquisition's output as a structural object (analytically neutral
    %       record). The form-vs-realization concern (delta computation,
    %       snapshot deletion) is handled by the form's role decomposition,
    %       not by this bullet.
    %   R3: sentence 2 tightened to remove repetition between "analytically
    %       neutral" and "no analytical commitment of its own". Verb-symmetry
    %       with sentence 1 — "with no commitments embedded by Acquisition"
    %       pairs with sentence 1's "embedded during capture".]
    \item \textbf{Independence from downstream interpretation.} At Acquisition's output, captured samples must carry no commitment to any specific downstream method, with no interpretation, transformation, or labeling embedded during capture. The captured sequence is therefore an analytically neutral record of what was captured and when, with no commitments embedded by Acquisition.
      % [src: structural paper §III.C.3]
  \end{itemize}

  % =========================================================================
  % SESSION PAUSED HERE — 2026-05-20 (resume tomorrow).
  %
  % §4.2.1 Acquisition polish status:
  %   - Purpose:        DONE (Variant 3, protocol + adjacency with gloss)
  %   - Responsibilities Bullet 1: DONE (R4 — unstructured-collection mix)
  %   - Responsibilities Bullet 2: DONE (R4 — scope built in, "order and
  %                                       temporal placement of captured states")
  %   - Responsibilities Bullet 3: DONE (R3 — verb-symmetry with sentence 1)
  %   - Invariant:      PROPOSED but not yet accepted. Three-sentence refined
  %                     version to consider on resume:
  %
  %     "At Acquisition's output, every sample in the observation sequence
  %      carries its temporal position, its adjacency relationship to
  %      neighbouring samples, and no commitment to any specific downstream
  %      method. If a downstream-specific transformation were embedded during
  %      capture (for example, structuring captured data into a particular
  %      analysis format or attaching labels from a specific method), the
  %      captured sequence would cease to be analytically neutral; what
  %      Acquisition emits would be tied to one method by construction. A
  %      realization satisfies this invariant when each emitted sample
  %      carries an explicit temporal position, an explicit relationship to
  %      its neighbouring samples, and no analytical commitment beyond the
  %      act of capture itself."
  %
  %   - Cross-reference: STILL FIRST-PASS. Needs polish in same style as
  %                      §4.1.1 / §4.1.2 Cross-references (light-formal-proof).
  %                      Likely Cross-reference type: structural handoff
  %                      (Acquisition's adjacency → Coordination's ordered
  %                      handoff).
  %
  % After §4.2.1: §4.2.2 Coordination polish, then re-read §4.2 as a whole,
  % then polish the §4.2 opener (apply the same artifacts-gloss treatment
  % that §4.1 opener got). Then §4.3 + sub-components.
  %
  % Pattern reminders from §4.2.1 polish (now also in updated SKILL.md):
  %   - Scope built into the claim ("At Acquisition's output", "during capture",
  %     "by Acquisition") avoids defensive "subsequent stages can..." punting.
  %   - "Memory evolution" carries semantic-loading risk; prefer structural
  %     phrasings like "order and temporal placement of captured states".
  %   - Sentence-2 patterns: "The result is X..." (Bullet 1, Bullet 2) /
  %     positive structural characterization of the output (Bullet 3).
  %   - Verb-symmetry between Sentence 1 and Sentence 2 within a bullet keeps
  %     repetition feeling structural rather than accidental.
  % =========================================================================
  \textbf{Invariant.} Each sampled state retains its position in the sequence and is collected without commitment to any specific downstream method.
  % [satisfies §3.1: comparability across the observed sequence (each sample
  %  occupies a known position in time); sampled states remain part of a
  %  distinguishable progression.
  %  Satisfies §3.2: observation should remain independent of downstream
  %  analysis assumptions.
  %  Satisfies §3.3: preserved separation between observation and downstream
  %  transformation.]

  \textbf{Cross-reference.} The state adjacency preserved here is what Coordination (\S4.2.2) carries forward as ordered handoff between stages. Both responsibilities derive from the comparability and ordered-relationships requirements in \S3.1.

  \subsubsection{Coordination}
  \textbf{Purpose.} Passes captured states between stages whose processing rates may differ, without losing their order or their link to the observational context in which they were acquired.

  \textbf{Responsibilities.}
  \begin{itemize}
    \item Decoupling of acquisition and transformation.
      % [src: structural paper §III.D.1 — Coordination Layer]
    \item Tolerance to variable processing latency.
      % [src: structural paper §III.D.2]
    \item Ordered handoff between stages.
      % [src: structural paper §III.D.3]
  \end{itemize}

  \textbf{Invariant.} The order of observations passed between stages is preserved regardless of the relative speeds at which those stages operate.
  % [satisfies §3.1: heterogeneous stage timing must not redefine the top-level
  %  ordering of the observed sequence.
  %  Satisfies §3.3: preserved separation between observation rhythm and
  %  transformation rhythm.]

  \textbf{Cross-reference.} The ordered handoff preserved here continues the state adjacency established by Acquisition (\S4.2.1). The decoupling of acquisition from transformation also marks the boundary between Component~2 and Component~3, ensuring that downstream transformation cannot redefine the order in which states were observed.

\subsection{Representation and Artifact Lifecycle}
% Component 3 opening.
% [first-pass — refine later]
Component~3 transforms observed state sequences into analysis-ready intermediate forms whose identity, interpretation, and persistence remain under explicit architectural control. It decomposes into three sub-components: \emph{Representation}, which produces the intermediate forms; \emph{Analysis}, which consumes them through methods that may evolve independently of the upstream roles; and \emph{Retention}, which governs whether artifacts persist while keeping persistence policy distinct from artifact meaning.

  \subsubsection{Representation}
  \textbf{Purpose.} Produces intermediate forms that make explicit the temporal relationship between consecutive observed states and that remain linked to the experimental step from which they originated.

  \textbf{Responsibilities.}
  \begin{itemize}
    \item Incremental representation of consecutive states.
      % [src: structural paper §III.E.1 — Representational Layer]
    \item Analysis-ready intermediate artifacts.
      % [src: structural paper §III.E.2]
    \item Step-scoped representation identity.
      % [src: structural paper §III.E.3]
  \end{itemize}

  \textbf{Invariant.} Each intermediate form preserves the temporal relationship of the states from which it was derived and remains traceable to its originating experimental step. Representation may \emph{organize} the observation sequence (for example, by aggregation or windowing) but must not \emph{redefine} it; the temporal order and identity of the underlying observations remain unaltered.
  % [Q2 resolved 2026-05-20: strengthened invariant to forbid "redefine" while
  %  permitting "organize" (aggregation, windowing). The user's distinction:
  %  organizing changes how the sequence is presented; redefining changes what
  %  it is.]
  % [satisfies §3.2: transformation into analysis-ready trace; preservation of
  %  identity of resulting outputs relative to the context from which they arise.
  %  Indirectly supports §3.1 by carrying temporal structure forward without
  %  redefining it.]

  \textbf{Cross-reference.} The step-scoped identity established here is the same identity that Retention (\S4.3.3) propagates through the artifact lifecycle. Both responsibilities derive from the artifact-identity requirement in \S3.2.
  % [Q3 resolved 2026-05-20: explicit cross-reference between Representation
  %  and Retention; both cite §3.2 as their common origin.]

  \subsubsection{Analysis}
  % [Q4 resolved 2026-05-20: Analysis kept as a peer sub-component, NOT folded
  %  into Representation. Rationale: §3.3 (Form-Level Requirements) demands
  %  that observation, transformation, representation, analysis, and artifact
  %  management remain distinct architectural separations. Collapsing Analysis
  %  into Representation would break that separation.]
  \textbf{Purpose.} Consumes the intermediate forms produced by Representation, through methods that may evolve or be deferred without forcing redesign of the upstream roles.

  \textbf{Responsibilities.}
  \begin{itemize}
    \item Downstream metric consumption.
      % [src: structural paper §III.F.1 — Analysis Layer]
    \item Deferred or modular evaluation.
      % [src: structural paper §III.F.2]
    \item Baseline-aware comparison support.
      % [src: structural paper §III.F.3]
  \end{itemize}

  \textbf{Invariant.} Analytical methods operate on intermediate forms, never on raw acquired states, so that they can be replaced or extended without disturbing upstream roles.
  % [satisfies §3.2: stable intermediate forms feeding later stages.
  %  Satisfies §3.3: form-level separation between roles — analysis remains
  %  one of the separations the form preserves.]

  \subsubsection{Retention}
  \textbf{Purpose.} Governs whether artifacts persist, are compressed, or are discarded, while keeping that policy distinct from the artifact's meaning in the experiment.

  \textbf{Responsibilities.}
  \begin{itemize}
    \item Configurable artifact persistence.
      % [src: structural paper §III.G.1 — Retention Layer]
    \item Separation of logical artifact value from storage policy.
      % [src: structural paper §III.G.2]
    \item Support for reproducible archival paths.
      % [src: structural paper §III.G.3]
  \end{itemize}

  \textbf{Invariant.} The persistence policy applied to an artifact is decided independently of the artifact's role in the experimental logic.
  % [satisfies §3.2: reproducible artifact lifecycle; traceable provenance.
  %  Satisfies §3.3: separation of practice from form — retention is part of
  %  the form's structure, but storage policy itself is outside the form's
  %  meaning, which is itself a form-level separation.]

  \textbf{Cross-reference.} The reproducible archival paths preserved here continue the step-scoped identity established by Representation (\S4.3.1). Both responsibilities derive from the artifact-identity requirement in \S3.2.
  % [Q3 resolved 2026-05-20: mirror of the Representation cross-reference;
  %  both sub-components explicitly trace the same identity chain back to §3.2.]

\subsection*{Mapping to Requirements (§3)}
% [Q1 resolved 2026-05-20: this Mapping subsection lives at the §4 wrapper
%  level (peer of §4.3, NOT inside it), so it can cover the full form once
%  Components 1 and 2 are drafted. This draft only develops Component 3,
%  so the §4.1 and §4.2 mapping paragraphs are placeholders below.]
% TODO: decide whether to render this as prose or as a table when the
% full form is drafted. Currently presented as structured paragraphs.

The Invariants of the architectural form satisfy the §3 requirements as follows.

% --- §4.1 Experimental Control mapping (drafted 2026-05-20)

\emph{Experimental Control} (\S4.1) addresses the requirements in \S3.1 that sampled states remain part of a distinguishable progression through the observed process and that the architecture support distinguishable observational conditions and units. Orchestration (\S4.1.1) preserves the step boundaries that make those conditions distinguishable without requiring semantic labels for what they represent; Stimulus (\S4.1.2) introduces the controlled, reproducible inputs that make changes in memory evolution relatable to changes in observed system activity. The explicit separation of Stimulus from Orchestration also satisfies the form-level requirement in \S3.3 that distinct architectural roles not collapse into one another.

% --- §4.2 Temporal Observation mapping (drafted 2026-05-20)

\emph{Temporal Observation} (\S4.2) addresses the requirements in \S3.1 that sampling preserve comparability across the observed sequence, that each sample occupy a known position in time, and that heterogeneous stage timing not redefine the top-level ordering of the observed sequence. Acquisition (\S4.2.1) captures successive states under a regime that preserves their adjacency and their positional identity, while remaining independent of downstream interpretation; Coordination (\S4.2.2) tolerates variable processing latency while preserving the order of states handed off between stages. The independence of Acquisition from downstream methods also satisfies the requirement in \S3.2 that observation remain independent of downstream analysis assumptions, and the form-level requirement in \S3.3 that observation and transformation remain distinct architectural roles.

% --- §4.3 Representation and Artifact Lifecycle mapping (drafted)

\emph{Representation} (\S4.3.1) addresses the requirement in \S3.2 that the observation sequence be transformed into an analysis-ready trace without altering the observational structure on which it depends, and the requirement that intermediate representations preserve their identity relative to the context from which they arise. It also indirectly supports the temporal-comparability requirement in \S3.1, by carrying the temporal relationships established upstream forward into the intermediate form rather than redefining them.

\emph{Analysis} (\S4.3.2) addresses the requirement in \S3.2 that stable intermediate forms be the basis for downstream stages, and the form-level requirement in \S3.3 that observation, transformation, representation, analysis, and artifact management not collapse into a single undifferentiated process. By forbidding analytical methods from operating on raw acquired states, the Invariant enforces exactly the separation between transformation and analysis that \S3.3 requires.

\emph{Retention} (\S4.3.3) addresses the requirement in \S3.2 that the architecture define a reproducible artifact lifecycle, specifying how artifacts are produced, transformed, preserved, reused, or discarded. By isolating persistence policy from the artifact's experimental meaning, the Invariant also satisfies the form-level requirement in \S3.3 that the form be characterized by the separations it preserves, including the separation between the form's meaning and engineering choices about how it is realized.

% NOTE: §3.1 (Temporal Observation Requirements) is not primarily owned by
% Component 3 — its primary owners are Components 1 (Experimental Control,
% for the conditions-and-units requirement) and 2 (Temporal Observation, for
% the sampling-and-comparability requirement). Component 3 only inherits
% §3.1 indirectly, by not destroying the temporal structure it receives.
% Cross-check: when Components 1 and 2 are drafted, verify they collectively
% own all of §3.1.



% =============================================================================
% NOTES
% =============================================================================
% -----------------------------------------------------------------------------
% §3 mapping cross-check (Step 3) — INITIAL CHECK (Component 3 only)
% -----------------------------------------------------------------------------
% I walked through §3.1, §3.2, §3.3 of paper-reordered.tex and verified
% coverage for Component 3.
%
% §3.1 Temporal Observation Requirements — coverage by Component 3:
%   - "sampled states must remain part of a distinguishable progression
%      ... each associated with its position in the sequence and the
%      observational condition under which it was acquired"
%        → primary owner: Components 1 and 2.
%        → Component 3 inherits: Representation Invariant carries the
%          temporal relationship of source states forward; Step-scoped
%          identity preserves the link to observational position.
%   - "sampling must also be defined in a way that preserves comparability
%      across the observed sequence"
%        → primary owner: Component 2 (Acquisition).
%        → Component 3 not responsible.
%   - "Temporal observation must also support distinguishable observational
%      conditions and units"
%        → primary owner: Components 1 (conditions) and 2 (units).
%        → Component 3 not responsible.
%   ⇒ §3.1 is NOT owned by Component 3. Flagged for Components 1/2 drafting.
%
% §3.2 Representation and Artifact Requirements — coverage by Component 3:
%   - "define how a volatile-memory observation sequence is transformed into
%      an analysis-ready trace without altering the observational structure"
%        → Representation Invariant (preserves temporal relationship of
%          source states) + Analysis Invariant (operates only on intermediate
%          forms, never on raw states).
%   - "observation and transformation must therefore remain distinct
%      architectural roles"
%        → enforced structurally by the separation between Component 2
%          (Observation) and Component 3 (Representation/Analysis/Retention).
%   - "preserve the identity of resulting outputs and intermediate
%      representations relative to the context from which they arise"
%        → Representation Invariant (traceable to originating step) +
%          Retention Responsibilities (reproducible archival paths).
%   - "define a reproducible artifact lifecycle, specifying how artifacts
%      are produced, transformed, preserved, reused, or discarded"
%        → Retention (preservation/discard) + Representation (production) +
%          Analysis (consumption/reuse).
%   ⇒ §3.2 is fully owned by Component 3.
%
% §3.3 Form-Level Requirements — coverage by Component 3:
%   - "defined at the level of form rather than as an inventory of mechanisms"
%        → satisfied by writing each sub-component in role-and-invariant terms.
%   - "defined by the separations it preserves... Observation, transformation,
%      representation, analysis, and artifact management should not collapse"
%        → satisfied by having Representation, Analysis, Retention as distinct
%          sub-components, each with its own Invariant; supplemented by
%          inter-component separation (§4.2 vs §4.3 when 4.2 is drafted).
%   - "must be characterized by its invariants"
%        → satisfied: every sub-component carries an explicit Invariant block.
%   - "general enough to admit multiple realizations, precise enough to guide
%      implementation and evaluation"
%        → realization-defending is deferred to §5 (Implementation) and
%          §6 (Analysis). The form itself is general (no mechanism mentioned).
%   ⇒ §3.3 is satisfied structurally by how Component 3 is presented.
%
% -----------------------------------------------------------------------------
% §3 mapping cross-check — FULL FORM (after §4.1, §4.2 drafted, 2026-05-20)
% -----------------------------------------------------------------------------
% Re-verified §3 coverage now that all three main components are drafted.
%
% §3.1 Temporal Observation Requirements — coverage by full form:
%   P1: "sampled states must remain part of a distinguishable progression
%        through the observed process, each associated with its position in
%        the sequence and the observational condition under which it was
%        acquired"
%        → §4.1.1 Orchestration (step boundaries → distinguishable progression);
%          §4.2.1 Acquisition Invariant (position in sequence);
%          §4.1.2 Stimulus + §4.1.1 Orchestration (observational condition).
%   P1: "preserve the phases or conditions under which states are sampled"
%        → §4.1.1 Orchestration; §4.1.2 Stimulus.
%   P2: "each sampled state should correspond to the same underlying state
%        abstraction; each sample should occupy a known position in time"
%        → §4.2.1 Acquisition (defined sampling regime; positional invariant).
%   P2: "Consistent sampling intervals... non-uniform sampling requires its
%        timing structure to be explicitly represented"
%        → §4.2.1 Acquisition (repeated state sampling responsibility).
%   P3: "distinguishable observational conditions and units"
%        → §4.1.2 Stimulus (conditions); §4.1.1 Orchestration + §4.2.1
%          Acquisition (units).
%   P3: "Differences in memory evolution... relatable to differences in
%        observed system activity without semantic knowledge"
%        → §4.1.2 Stimulus (reproducible workload, varied across runs);
%          §4.2.1 Acquisition (semantic-neutral, independent of downstream).
%   P3: "distinguishable segments... preserve segments without semantic labels"
%        → §4.1.1 Orchestration (boundaries); §4.2.2 Coordination (ordered
%          handoff carries segment structure forward).
%   ⇒ §3.1 is now fully owned by Components 1 and 2.
%     The flag from the initial Component-3-only check is RESOLVED.
%
% §3.2 Representation and Artifact Requirements — coverage by full form:
%   - "transformed into an analysis-ready trace without altering the
%      observational structure" → §4.3.1 Representation Invariant
%      (organize but not redefine).
%   - "Observation and transformation must therefore remain distinct
%      architectural roles" → boundary between §4.2 and §4.3, reinforced
%      by §4.2.1 Acquisition (independence from downstream interpretation)
%      and §4.3.2 Analysis (operates only on intermediate forms).
%   - "observation should remain independent of downstream analysis
%      assumptions" → §4.2.1 Acquisition Invariant (now explicitly).
%   - "preserve the identity of resulting outputs and intermediate
%      representations relative to the context from which they arise"
%      → §4.3.1 Representation (step-scoped identity); §4.3.3 Retention
%      (reproducible archival paths).
%   - "reproducible artifact lifecycle: produced, transformed, preserved,
%      reused, or discarded" → §4.3 Component 3 collectively.
%   ⇒ §3.2 is fully owned. Component 3 is the primary owner, with §4.2
%     contributing the observation-independence half of the separation.
%
% §3.3 Form-Level Requirements — coverage by full form:
%   - "defined at the level of form rather than as an inventory of
%      mechanisms" → all sub-components written in role-and-invariant terms.
%   - "defined by the separations it preserves" — explicit separations now
%     enumerated at three levels:
%       (a) intra-component: Orchestration vs Stimulus (§4.1);
%           Acquisition vs Coordination (§4.2);
%           Representation vs Analysis vs Retention (§4.3).
%       (b) inter-component: §4.1 vs §4.2 vs §4.3.
%       (c) cross-cutting: storage policy vs artifact meaning (§4.3.3),
%           analysis methods vs raw states (§4.3.2).
%   - "characterized by its invariants" → seven sub-components, each with
%     one explicit Invariant.
%   - "general enough to admit multiple realizations, precise enough to
%      guide implementation and evaluation" → realization-defending is
%     deferred to §5 (Implementation) and §6 (Analysis).
%   ⇒ §3.3 is fully satisfied by the structural pattern of the form.
%
% FULL-FORM VERDICT: every §3 requirement now has at least one owner in
% §4.1, §4.2, or §4.3. The form is internally complete with respect to §3.
%
% -----------------------------------------------------------------------------
% Resolved questions (decisions logged 2026-05-20)
% -----------------------------------------------------------------------------
%   Q1. Mapping subsection placement.    [RESOLVED — move up]
%       Decision: the Mapping subsection lives at the §4 wrapper level so it
%       can cover the full form, not just Component 3. Placeholders for §4.1
%       (Experimental Control) and §4.2 (Temporal Observation) added to the
%       Mapping subsection above; they will be filled when those components
%       are drafted.
%
%   Q2. Representation Invariant breadth.   [RESOLVED — strengthen with nuance]
%       Decision: the Representation Invariant explicitly forbids redefining
%       the observation sequence, but explicitly permits ORGANIZING it (for
%       example, aggregation or windowing). The user's distinction: organizing
%       changes how the sequence is presented; redefining changes what it is.
%       The strengthened invariant is in §4.3.1.
%
%   Q3. Cross-component coupling between Representation and Retention.
%                                            [RESOLVED — explicit cross-refs]
%       Decision: both sub-components carry an explicit Cross-reference
%       block. §4.3.1 (Representation) points to §4.3.3 (Retention) for
%       lifecycle continuation; §4.3.3 (Retention) points back to §4.3.1
%       (Representation) for identity origin. Both cite §3.2 as their
%       common requirement source.
%
%   Q4. Whether Analysis deserves its own sub-component.
%                                            [RESOLVED — keep as peer]
%       Decision: Analysis stays as a peer sub-component of Representation
%       and Retention. Collapsing it into Representation would break the
%       §3.3 (Form-Level) requirement that observation, transformation,
%       representation, analysis, and artifact management remain distinct
%       architectural separations. Rationale logged as a comment in §4.3.2.
%
% -----------------------------------------------------------------------------
% Items NOT changed (per user instruction)
% -----------------------------------------------------------------------------
%   - paper-reordered.tex is unchanged. The §4 body in that file remains
%     the TODO placeholder. This draft is meant to be reviewed and then
%     copied or \input{}'d into paper-reordered.tex once accepted.
%
% =============================================================================
